\section{Introduction}
Concept erasure edits a text-to-image model so that it stops producing a target
concept, such as an artist's style or nudity. Erased checkpoints are audited in full
precision, but they are often deployed after post-training quantization. The audited
object and the deployed program therefore differ. Two failure modes follow. Passively,
quantization may round a small edit back toward the original weights, as reported for
LLM unlearning at 4 bits~\cite{zhang2025}. Adversarially, a supplier may ship weights
that look erased in full precision but deploy as the unedited model, extending
quantization-conditioned attacks on LLMs~\cite{egashira2024,egashira2025} to erasure.

We study both on Stable Diffusion~1.5 with Unified Concept Editing
(UCE)~\cite{uce}, which edits only the 32 cross-attention key and value
projections. This restriction is useful rather than limiting. The inputs to these
projections are text-encoder embeddings that do not depend on the UNet, so the edit's
entire behavioral effect passes through a linear map that we can analyze exactly.
Our contributions are:

\begin{itemize}
\item \textbf{A characterization of deployed collisions.} Scale-preserving shipped
weights that deploy like the original form a union of boxes indexed by the coordinates
attaining row and column maxima (Prop.~\ref{prop:box}). Pinning the original arg-max
pattern yields one box with both row and column caps. This constraint explains why an
earlier iterative construction stalled at 83--84\% code agreement; the new constructed
checkpoints attain exact agreement.
\item \textbf{Optimal reconstruction in a specified box.} Because images depend on
the edited weights only through their key/value outputs, we solve a convex program
for the best reconstruction of this edit within the original-pattern box
(Eq.~\ref{eq:capacity}). Its duality gap certifies reconstruction in that box on the fit
prompts, not optimal behavioral stealth or the full union. We compute it for W8A8,
INT8, INT4, and NF4.
\item \textbf{A paired behavioral evaluation.} Using identical prompts and seeds across
FP16, simulated W8A8, native INT8, and 4-bit paths, with a crossed seed--prompt
bootstrap, we measure honest deployment gaps for style and nudity erasure and evaluate
the optimal attack with a stealth-first criterion.
\item \textbf{An image-free diagnostic.} Conditional on key/value reconstruction,
we bound a deployment-change statistic (Prop.~\ref{prop:audit}) and measure how it
separates the evaluated attacks from honest checkpoints without generating images.
\end{itemize}

The results split cleanly by bit width. At 8 bits, with or without SmoothQuant, only
16\% of edited coordinates fit in the selected box. Its reconstruction optimum has
a certified fit-prompt residual of at least 0.64, while the constructed checkpoint
reproduces 34--42\% of the edit's measured effect on held-out prompts and scores
like the unedited model in FP16. This is not a lower bound on the behavioral stealth
of every feasible checkpoint. At 4 bits the constructed checkpoints show recovery on
both concepts. The INT4 nudity checkpoint has an FP16 point score near the honest
erasure (NudeNet
0.168 against 0.173) but deploys bit-for-bit as the 4-bit original (0.595), a
recovery of $+0.43$. The FP16 similarity is descriptive rather than established
equivalence. Honest erasures, in contrast, show no significant deployment
gap under any tested recipe, including whole-UNet W8A8, INT4, and NF4, and benign
prompt alignment is unchanged. The one warning sign is whole-UNet INT4, which raises
the honest nudity erasure's score by a prompt-dependent $+0.11$ that our 64-image
control cannot resolve. The proposed diagnostic separates the evaluated 4-bit
attacks at 2.2--3.3 times the honest value, without a calibrated false-positive rate.
